ThinkStation PGX% =====================================================================
% BID2---MOTIVATION SECTION (standalone; to be placed into the paper)
% =====================================================================
% Build date:      2026-08-09
% Built by:        Claude Code (Fable 5), directed by Héctor J. Villarreal
% Build instruction: motivation/CC_instrucciones_motivation-section-draft_v1.md
% Branch:          p3-correcciones-tex
%
% NUMBER SOURCE. Every numeral in the prose below is taken from
%   motivation/output/NUMBERS.md, motivation/output/VERIFICATION.md, or the
%   per-figure tables in motivation/output/tables/ (read from the CSV, never
%   off a figure). All ENIGH money figures are quarterly pesos per person or
%   per household in constant Aug–Nov 2024 pesos (general INPC).
%
% PLACEMENT (suggestion only---Héctor and Diego decide). This reads as a new
%   section immediately following the introduction (\ref{sec:intro}) and
%   before the literature review, or alternatively as the opening of
%   Section~\ref{sec:data}. It does not renumber or modify Draft-June-v6.tex.
%   Labels use the prefix "mot_" and collide with nothing in that file.
%
% CITATION ROUTE. references.bib (BID2 root, 18 entries) has NO keys for
%   ENIGH, ENSANUT, ENASEM, GHED, Searle et al. (2008) or Hosseini et al.
%   (2021). Following the convention already used in Draft-June-v6.tex ---
%   where ENASEM, ENIGH, ENOE and the WHO life tables are named in prose with
%   "Authors' calculations" and no \cite---these are written as plain text so
%   the section compiles against the existing bibliography with no .bib edits.
%   If Héctor prefers proper entries, the six keys can be added later without
%   touching the prose beyond wrapping the names in \citep{}.
%
% DECISION FLAGGED FOR HÉCTOR---fig 10 survival panel. The mortality panel of
%   fig10_enasem_dynamics is NOT age-adjusted, and unconditional mortality by
%   wealth tercile is flat to inverted (5.1 / 6.7 / 6.6 percent,
%   VERIFICATION.md §5)---almost certainly age composition, since wealth
%   accumulates and higher terciles are older in a 50+ panel. The conditional
%   gradient is also non-monotone (17.4 / 20.1 / 10.5 percent). DECISION TAKEN
%   HERE: the survival gradient is OMITTED from the body; the spending-response
%   result carries the ENASEM evidence, which is the finding the mechanism
%   needs. If Héctor wants it reinstated, the drafted sentence is parked as a
%   comment at the end of §\ref{subsec:mot_unmet}. The age adjustment is a
%   separate scripted rerun, not attempted here.
%
% DISCREPANCY NOTED (not reproduced in the prose). READTHIS.md and the build
%   instruction paraphrase catastrophic incidence as falling from "~4–5% in
%   the poorest deciles to about 1% in the richest". The tables do not support
%   the upper end: in 2024 the WHO 40%-of-capacity-to-pay measure is 3.90% in
%   decile 1 and 1.37–1.88% in deciles 4–10 (1.83% in decile 10); in 2018 it is
%   2.49% and 1.40%. The prose below quotes the table values.
%
% FIGURE MANIFEST---see end of file.
% =====================================================================

\section{Health Spending in Mexican Micro-Data}
\label{sec:motivation}

The model developed below places health at the analytical centre, and it does so
because of what household data say about how Mexican families actually spend on
medical care. This section establishes four empirical premises from three
independent surveys: that medical spending has a steep and stable age--sex
profile; that the capacity to finance it is deeply unequal; that the spending is
majority curative rather than preventive; and that low spending among poor
households reflects unmet need rather than good health. Two conventions apply
throughout. All spending measured here is \emph{out-of-pocket}: the in-kind
component delivered through IMSS, ISSSTE and the SSA network does not appear
anywhere in these profiles, so every level below understates the resources
actually devoted to a household's health. And the household survey used for the
profiles, ENIGH, captures only 22 to 34 percent of the out-of-pocket per-capita
spending recorded for M\'exico in the WHO Global Health Expenditure Database;
its evidence is therefore used for shapes, gradients and ratios, never for
levels or national totals.

\subsection{An age--sex profile that is not an artifact}
\label{subsec:mot_profile}

Out-of-pocket medical spending rises monotonically with age from young adulthood
onward. Under a utilization-weighted allocation of household spending to
individuals, a woman aged 85 or over is assigned roughly five times what a woman
aged 15--29 is assigned, and a man in the oldest group roughly six and a half
times his young-adult counterpart (Figure~\ref{fig:mot_profile}). Because
individual spending is not observed directly in a household survey, the natural
objection is that this profile is manufactured by the allocation rule. It is
not. Two methods that share no assumptions---a household-level regression of
total health spending on age--sex composition, and an NTA-style allocation using
utilization weights estimated from ENSANUT---produce profiles that correlate at
0.959 across the eighteen age--sex cells of the 2024 wave. Where the two are
compared cell by cell they agree closely at exactly the ages that matter: for
women aged 65--74 the regression gives 1{,}149 quarterly pesos against 1{,}241
from the allocation, and for women aged 75--84, 1{,}614 against 1{,}501. The
shape is also stable across the four ENIGH waves from 2018 to 2024. The profile
is an estimate, and it is labelled as such wherever it is used, but it is not an
artifact of one arbitrary rule.

Two features of the profile matter for what follows. First, women spend
distinctly more than men through the reproductive and midlife years---about 32
percent more at ages 30--44 and 38 percent more at ages 45--59---and the gap
narrows and can reverse at the oldest ages. Second, the estimated profile can be
checked against a survey that observes individual spending directly. ENASEM, the
Mexican longitudinal study of ageing, yields a 50-and-over profile that rises
with age in the same way but sits three to four times higher in level: 3{,}836
quarterly pesos for men aged 75--79 in 2024 against the 1{,}099 the ENIGH-based
allocation assigns to men aged 75--84. Two independent comparisons, one
external and one internal, therefore point the same way, and they are the reason
this paper treats the profile as information about shape alone.

\subsection{The capacity to spend is unequal}
\label{subsec:mot_inequality}

Health's share of the household budget looks almost flat across the income
distribution: 3.5 percent in the poorest decile and 3.7 percent in the richest in
2024, with a shallow trough in between. That flatness is an artifact of
zeros. Among households with no member aged 65 or over, 45.5 percent of the
poorest decile reported no health spending at all in the reference quarter,
against 27.9 percent of the richest; the gap barely closes when the household
does contain someone 65 or over (39.5 against 24.8 percent). Conditional on
spending anything, the poorest decile devotes 6.1 percent of its budget to
health, more than any other decile and well above the 3.9 to 5.1 percent
observed higher up. The poor do not spend less because health absorbs a smaller
share of their resources; they spend less because they more often spend nothing.

When they do spend, the burden is heavier. Catastrophic spending on the WHO
capacity-to-pay definition---health exceeding 40 percent of resources net of
subsistence---reaches 3.9 percent of households in the poorest decile in 2024
against 1.4 to 1.9 percent in deciles four through ten, and the ordering is
unchanged under the alternative threshold of 25 percent of total expenditure.

The premise the model leans on most heavily, however, is that inequality lives
\emph{inside} the age profile rather than beside it. Re-estimating the profile
separately by income tercile, a woman aged 75--84 in the top tercile is assigned
2{,}545 quarterly pesos against 782 in the bottom tercile, a ratio of 3.3. The
same comparison at ages 15--29 gives 446 against 150. In absolute terms the
climb from young adulthood to old age is 2{,}099 pesos in the top tercile and
632 in the bottom. Ageing and inequality are not two separate forces acting on
the household budget; the second determines how steeply the first bites.

\subsection{Spending responds to illness rather than preventing it}
\label{subsec:mot_curative}

Every expenditure item in the ENIGH health module was classified as
unambiguously preventive, unambiguously curative, or genuinely ambiguous, under
a stated rule recorded item by item in the build outputs. This is an auditable
judgment, not a measured share, and it is reported with the bracket that
judgment implies. In 2024, unambiguously preventive items account for 2.0
percent of out-of-pocket health spending---a figure between 1.9 and 2.5 percent
in every wave---unambiguously curative items for 53.6 percent, and the remaining
44.5 percent is ambiguous: over-the-counter medicines, dental and optical
services, chronic-disease maintenance, unqualified diagnostics and insurance
premiums (Figure~\ref{fig:mot_composition}). Assigning the entire ambiguous
block to prevention still leaves curative spending the majority, at 53.6 against
46.4 percent; assigning it to treatment makes spending 98.1 percent
curative. The licensed conclusion is therefore the weaker and more robust one:
Mexican out-of-pocket health spending is \emph{majority curative}, and the
preventive share lies somewhere between 2 and 46 percent depending on where the
ambiguous block is assigned. Because the 2024 COICOP recode of the expenditure
classification dropped prescription status, this composition is quoted
cross-sectionally only and no trend across waves is claimed.

A second survey observes the motive of care directly rather than inferring it
from the receipt, and it agrees. In ENSANUT 2024, preventive motives---vaccination, check-up, prenatal control---account for 11.2 percent of the last
health need reported in the poorest socioeconomic tercile, rising only to 15.9
percent in the richest. Prevention is a minority of care everywhere and least
common where resources are scarcest.

\subsection{Low spending among the poor is unmet need}
\label{subsec:mot_unmet}

If poor households spend little because they are healthy, the zeros documented
above are benign and there is nothing for a model of health to explain. Three
independent lines of evidence say they are not. The first is the pattern of
zeros itself: they are concentrated among poor households and remain elevated
there even when the household contains an older member, which is precisely where
need is least plausibly absent. The second is direct measurement. Among people
who reported a health need in ENSANUT 2024, 19.2 percent in the poorest
wellbeing tercile went without care, against 8.9 percent in the richest; forgone
care attributed specifically to cost collapses across the same gradient, from
3.5 percent to 0.1 percent (Figure~\ref{fig:mot_forgone}). The uninsured are the
worst off on both margins. The 2018 wave, using the survey's coarser
socioeconomic stratum, shows the same shape. Cost-related unmet need is a
measured fact in these data, not an inference from a spending gradient.

The third line is dynamic. Following individuals in ENASEM whose health deficits
increased materially between 2018 and 2021 (a rise of at least 0.10 in the
index, about a decile of its range), out-of-pocket spending rises in every
wealth tercile---but the increase is associated with wealth, roughly doubling
from 10{,}819 annual pesos in the bottom tercile to 21{,}230 in the top, on
2021 spending levels of 24{,}043 against 42{,}233. The same shock is met with
twice the response at the top of the wealth distribution. This pattern is
descriptive and not causal: wealth in these data bundles insurance, education,
prior health status and access to care, and no attempt is made here to separate
them. The frailty index used to define the health decline is a reconstruction of
the deficit-accumulation index of Searle et al.\ (2008), following the
convention used in the project's earlier ENASEM work and reproducing its
published results qualitatively rather than numerically; it is provisional
pending confirmation, and the pattern reported here does not hinge on the fine
details of its construction.
%
% PARKED FOR HÉCTOR (fig 10 right panel; see header). If the survival gradient
% is to be reinstated, insert after the spending-response sentence:
%   "Conditional on a comparable decline in health, mortality by 2024 is 10.5
%   percent in the top wealth tercile against 17.4 and 20.1 percent in the lower
%   two; this comparison is not age-adjusted, the terciles differ in age
%   composition, and the ordering across the lower two terciles is not
%   monotone, so it is reported only as suggestive."

\subsection{The 2020 episode}
\label{subsec:mot_pandemic}

The pandemic wave of ENIGH converts the cross-sectional reading of the zeros
into an observation in time. Between 2018 and 2020 real household out-of-pocket
health spending rose in every income decile, by 76 percent in the poorest and 29
percent in the richest, and the rise never reversed through the 2022 and 2024
waves (Figure~\ref{fig:mot_pandemic}). The margin of adjustment was
participation rather than intensity: the share of households spending nothing on
health fell by 13 to 15 percentage points in every decile, while spending
conditional on being positive rose far less---36 percent in the poorest decile
and 7.5 percent in the richest. A placebo on food spending shows nothing
comparable: food fell in most deciles and rose 10 percent in the poorest, where
total monetary expenditure also rose about 10 percent, against health's 76.

Two framings of the same episode point in opposite directions and both are
reported here. In growth rates the shock is regressive, hitting the poorest
hardest. In pesos it is not: the top two deciles still financed 36 percent of
the aggregate increase. What the participation margin shows is that when the
free public option became unavailable, poor households were conscripted into
paying out of pocket rather than priced out of responding---which is only
possible if the need had been there all along, absorbed by in-kind provision.
The causal language stays deliberately soft: two waves and one placebo cannot
fully exclude a health-specific change in reporting, though the same instrument
was used in both waves. The episode merits separate treatment, which this paper
does not pursue.

\subsection{What the model must contain}
\label{subsec:mot_bridge}

Taken together these facts discipline the model that follows. Medical spending
is a response to a realized health shock rather than a smooth investment in a
health stock, so health must enter as a stock subject to shocks and medical
spending as restoration rather than accumulation. The capacity to finance that
restoration is unequal, and the inequality is embedded in the age profile
itself, so the budget constraint must bind differently by type and must do so
more tightly with age. And because the constraint binds, a transitory health
shock is not transitory in its consequences: the household that cannot finance
restoration carries the deficit forward. The model in Section~\ref{sec:model}
is built to reproduce exactly these features.

% ---------------------------------------------------------------------
% FIGURES
% ---------------------------------------------------------------------

\begin{figure}[H]
\centering
\includegraphics[width=0.95\textwidth]{fig01_age_sex_profile.pdf}
\caption{Estimated out-of-pocket medical spending by age and sex, M\'exico 2024.
Quarterly pesos per person, constant Aug--Nov 2024 pesos (general INPC).
Individual spending is not observed in ENIGH; the profile is \emph{estimated}
under two independent allocation rules---a household regression on age--sex
composition and an NTA-style allocation with utilization weights from ENSANUT
2024---which correlate at 0.959 across the eighteen age--sex cells. Design-based
estimates (strata $\times$ PSU $\times$ survey weights). Out-of-pocket spending
only; in-kind provision through IMSS, ISSSTE and SSA is excluded. Source:
ENIGH 2024 (INEGI) and ENSANUT Continua 2024 (INSP), authors' calculations.}
\label{fig:mot_profile}
\end{figure}

\begin{figure}[H]
\centering
\includegraphics[width=0.95\textwidth]{fig07_composition.pdf}
\caption{Composition of out-of-pocket health spending by income decile under the
stated curative/preventive classification. Every expenditure item is assigned to
one of three classes by an auditable rule; the ambiguous block (over-the-counter
medicines, dental, optical, chronic-disease maintenance, unqualified diagnostics,
insurance premiums) is reported rather than allocated. Assigning it entirely to
prevention leaves curative spending the majority; assigning it to treatment makes
spending 98.1 percent curative. Because the 2024 COICOP recode dropped
prescription status, the composition is comparable across deciles within a wave
but not across waves. Source: ENIGH 2024 (INEGI), authors' calculations.}
\label{fig:mot_composition}
\end{figure}

\begin{figure}[H]
\centering
\includegraphics[width=0.95\textwidth]{fig06_forgone_care.pdf}
\caption{Forgone care among people reporting a health need, by socioeconomic
tercile and by insurance affiliation, ENSANUT Continua 2024. The lighter series
restricts forgone care to needs the respondent attributed to cost. Design-based
estimates. Source: ENSANUT Continua 2024 (INSP), authors' calculations.}
\label{fig:mot_forgone}
\end{figure}

\begin{figure}[H]
\centering
\includegraphics[width=0.95\textwidth]{fig12_pandemic_jump.pdf}
\caption{Change in real household out-of-pocket health spending by income
decile, 2018 to 2020, with food spending as a placebo and the change in the
share of households spending zero. Constant Aug--Nov 2024 pesos (general INPC).
The 2020 wave is pandemic-affected; the disruption is the object of interest
here rather than a nuisance. Source: ENIGH 2018 and 2020 (INEGI), authors'
calculations.}
\label{fig:mot_pandemic}
\end{figure}

% ---------------------------------------------------------------------
% METHODS FOOTNOTE---attach to the first sentence of §\ref{subsec:mot_profile}
% if Héctor wants the machinery visible in the section rather than the appendix.
% ---------------------------------------------------------------------
% \footnote{All estimates are design-based (ENIGH \texttt{est\_dis} strata,
% \texttt{upm} PSUs, \texttt{factor} weights; analogous designs for ENSANUT and
% ENASEM). ENIGH health items are selected from each wave's shipped expenditure
% catalogue, not from memory, and reproduce the published \texttt{salud}
% aggregate to within 0.4 percent in every wave. Money is deflated to Aug--Nov
% 2024 pesos with the general INPC. Full item lists, allocation code and every
% plotted value are in the build outputs.}

% ---------------------------------------------------------------------
% ADDITIONAL EVIDENCE AVAILABLE (referenced in prose, not reproduced here).
% Available for an appendix if wanted: validation of the estimated profile
% against ENASEM's directly observed 50+ profile; the profile by income tercile;
% zero-spending by decile and 65+ status; catastrophic incidence by decile under
% both thresholds; the ENASEM spending-response and outcome panels; and the
% four-wave stability of the profile shape.
% ---------------------------------------------------------------------

% =====================================================================
% FIGURE MANIFEST---files to upload to the Overleaf project root
% =====================================================================
%   fig01_age_sex_profile.pdf
%   fig07_composition.pdf
%   fig06_forgone_care.pdf
%   fig12_pandemic_jump.pdf
%
% Source: Missions/Funded/BID2/motivation/output/figures/
% (.png versions of each exist in the same folder if PDF inclusion is a problem;
% Draft-June-v6.tex currently uses .png throughout.)
%
% If an appendix of supporting figures is added, the additional files are:
%   fig02_validation_enigh_vs_enasem.pdf, fig03_profile_by_income.pdf,
%   fig04_budget_share_by_decile.pdf,     fig05_zero_spending.pdf,
%   fig06b_forgone_by_insurance.pdf,      fig08_catastrophic.pdf,
%   fig09_spending_distribution.pdf,      fig10_enasem_dynamics.pdf,
%   fig11_wave_stability.pdf
% =====================================================================
